% Study Characteristics Subsection
\subsection{Study Characteristics}
\label{sec:study-characteristics}

The 13 studies included in this review all vary in sample size and study design to a great degree. Detection studies enrolled a median of 44 participants (range 1 to 384), while forecasting studies enrolled a median of 11 participants (range 6 to 70).

Three studies employed prospective designs. \textcite{spahrDeepLearningbasedDetection2025} conducted a multi-center prospective study across eight EMUs. \textcite{dongTwoLayerEnsembleMethod2022} prospectively monitored patients at home for up to three months. \textcite{fineDetectionKeyAutomated2025} prospectively recorded data from 15 patients in an EMU for tonic seizure detection. The remaining studies used retrospective analysis of existing datasets or small-scale feasibility designs.

Research activity on DL models for detection has accelerated over time. Three forecasting studies were published between 2020 and 2022. Nine detection studies were published between 2023 and 2026. 

Seven studies were conducted in hospital or EMU settings. Five studies included home or ambulatory validation and \textcite{elemamAutomatedValidatedTool2025} used nighttime recordings captured at homes and in-hospital.

Four studies used a sample size of less than 20 participants. The smallest sample size was used by \textcite{singhrathoreDevelopmentMultimodalMachine2024}. This is  a single-patient case study. Also \textcite{borujenyDetectionEpilepticSeizure2013} evaluated only 3 patients. Small sample sizes limit the reliability of performance estimates and generalizability.
